\documentclass[11pt]{article}
\usepackage{graphicx}
\usepackage[a4paper,margin=.5in]{geometry}
\usepackage{amsmath}
\usepackage{booktabs}
\usepackage{array}
\usepackage{multirow}
\usepackage{pgfgantt}

\usepackage[
backend=biber,
style=apa,
]{biblatex}
\addbibresource{./bib.bib}
\usepackage{helvet}
\renewcommand{\familydefault}{\sfdefault}

\title{\textbf{Project Brief: Lifetime Asset Allocation in New Zealand}}
\author{Callum Davidson}
\date{\today}

\begin{document}

\maketitle

\section{Domain}

The research addresses the critical question of how individuals should allocate their retirement savings across different asset classes throughout their lifetime, both during the accumulation phase (working years) and the decumulation phase (retirement). This project operates at the intersection of quantitative finance, retirement planning, and public policy in the New Zealand context.

The domain presents some unique challenges. First, the mathematical complexity of lifetime allocation problems makes them difficult to solve analytically, particularly when incorporating realistic investor behaviour such as fixed-dollar withdrawals rather than proportional spending. Second, New Zealand's retirement landscape differs significantly from the U.S.-centric academic literature. New Zealand Superannuation provides a universal pension regardless of means, unlike comparable international systems. The UK combines state pensions with means-tested pension credits \parencite{ukgov2025pensionCredit,ukgovNDpensionExplained}, while US Social Security benefits are earnings-linked \parencite{ssa2025beniftsPDF} and Australia's Age Pension is means-tested \parencite{ausgov2025agePension}. These policy differences significantly influence retiree behaviour and risk tolerance, limiting the generalisability of international research.

Moreover, KiwiSaver, New Zealand's workplace-based retirement savings scheme, has unique features such as government contributions and employer matching that affect accumulation dynamics. Australia has a similar system, but New Zealand's differs in that it has lower contribution rates and the ability for New Zealanders to withdraw funds for first-home purchases. In Australia, you can only withdraw up to 15,000 AUD for first-home purchases, whereas in New Zealand, you can withdraw your entire KiwiSaver balance (subject to certain conditions) \parencite{firstsuper2025kiwisaverVsSuper}.

The significance of this research extends beyond individual welfare. The global asset-backed pension market was valued at 68.9 trillion USD in 2024 \parencite{oecd2025pension}. In New Zealand, demographic pressures on the retirement system are mounting, with Treasury projecting that spending on New Zealand Superannuation and healthcare could account for nearly half of all government expenditure by 2065 \parencite{nztreasury2025tirohanga}. There is broad consensus that New Zealanders are under-saving for retirement \parencite{national2025kiwisaver, gilbert2025lifting, palmer2025watch, dexter2025kiwisaver}. Lifetime asset allocation strategies that improve retirement outcomes could help alleviate some of these pressures by enhancing the efficiency of individual savings and reducing reliance on government support. However, there are significant tensions between industry practice (which predominantly recommends declining-risk strategies during decumulation) and academic research (which often finds constant or even increasing-risk strategies to be superior in decumulation) \parencite[e.g.][]{estrada2014glidepath, estrada2015retirement, anarkulova2025beyond}. The lack of consensus creates confusion for both individual investors and financial advisers, potentially leading to suboptimal retirement outcomes. This research aims to bridge the gap between theory and practice by developing a flexible simulation and optimisation framework tailored to the New Zealand context, evaluating existing glide-path recommendations, and providing evidence-based guidance for investors and policymakers.

This research is conducted in partnership with Consilium NZ LTD, a New Zealand-based financial services company with whom the author is currently employed as a Data Engineer. Consilium provides investment advisory services and retirement planning solutions to individual and institutional clients. The partnership provides access to proprietary market data and industry insights that will enhance the relevance and applicability of the research findings. However, the author's academic independence is maintained, with the author retaining intellectual property rights and research autonomy. As such, no confidential or proprietary information from Consilium is used in this research.


\section{Goal}

This research has four primary objectives. The first is to develop a flexible simulation and optimisation framework for lifetime asset allocation that addresses key limitations in existing approaches. Many methods rely on restrictive assumptions that limit their practical applicability. For instance, the foundational work of \textcite{merton1969lifetime} assumed the investor made fixed withdrawals of a proportion of current wealth, which does not reflect how most retirees actually spend. Similarly, many studies assume that asset returns are independent and identically distributed, an assumption that \textcite{anarkulova2025beyond} have argued leads to misleading conclusions about the benefits of bonds in retirement portfolios. Our framework should therefore remain agnostic to the source of asset returns and allow for realistic spending patterns that reflect actual retiree behaviour (variable dollar withdrawals rather than fixed percentage withdrawal).

The second objective is to evaluate existing glide-path recommendations within the New Zealand context using the framework developed. We wish to pay particular attention to testing the declining-risk strategies that dominate industry practice. Target-date funds in the United States predominantly follow declining-risk structures \parencite{dolvin2010asset}, and similar patterns exist in New Zealand products such as SuperLife's Age Steps \parencite{superlife_nodate}. The intuitive appeal of these strategies is straightforward: older investors have less time to recover from losses, making a gradual shift towards safer assets seem sensible. However, the academic literature does not strongly support this approach, particularly during the decumulation phase \parencite[e.g.][]{estrada2014glidepath, estrada2015retirement, anarkulova2025beyond,shiller2005life,merton1969lifetime}. By applying our simulation framework with realistic assumptions about New Zealand retirees, including appropriate spending patterns and risk concerns, we can assess whether the theoretical arguments against declining-risk strategies hold in practice. This must necessarily account for New Zealand policy settings, particularly the impacts of New Zealand Superannuation, KiwiSaver, and the Healthcare System. 

The third objective involves extending existing optimisation methods to find lifecycle asset allocation strategies that meet the needs and concerns of real-world investors. This will involve researching and identifying suitable risk or utility measures that capture retirees' practical concerns, such as the probability of portfolio depletion at different ages or total consumption shortfall, rather than relying solely on terminal wealth variance or expected utility. We will then implement optimisation methods to identify asset allocation strategies that vary based on the state of the investor, including both age and wealth levels. This objective comprises two distinct sub-goals. First, we will replicate the findings of \textcite{anarkulova2025beyond} regarding age-based glide-paths, but under spending assumptions that better reflect New Zealand retiree behaviour. This replication serves both to validate our framework against published results and to examine how their conclusions change under different spending scenarios. Second, we will implement the control-matrix approach proposed by \textcite{Makinen2024}, which allows asset allocation to vary with both time and wealth. However, their original formulation has some limitations that we will seek to address. Their risk measure focuses exclusively on terminal wealth, which fails to capture the holistic, whole-lifecycle nature of risk that concerns real investors. Additionally, we also believe that the way that the control matrix is constructed could be improved to enhance optimisation efficiency. The optimisation process on such a high-dimensional parameter space will be the key technical challenge. We will explore gradient-based optimisation methods, leveraging automatic differentiation tools available in frameworks like PyTorch, as well as any other suitable optimisation algorithms identified during the literature review phase and consultations with academic experts.

The fourth objective is to provide evidence-based recommendations for individual investors and financial advisers in New Zealand. We will evaluate Consilium's existing retirement products, identifying opportunities for improvement. This includes critical assessment of our findings' limitations and ethical communication considerations. Given inherent market uncertainty, recommendations must acknowledge both analytical strengths and limitations in language accessible to advisers and clients, avoiding overconfidence in any particular strategy.

If time constraints prevent achieving all objectives, priority will be given to the first two. The first is essential as it creates the analytical infrastructure for all subsequent work. The second has the most immediate practical relevance, directly addressing whether current industry practices are evidence-supported. The third objective represents an important methodological contribution, but the control-matrix extension may be deferred if earlier phases require more time. The fourth objective will be integrated throughout, with preliminary recommendations developed as findings emerge.


\section{Data}

This research requires two primary categories of data: asset return data to drive Monte Carlo simulations, and retirement spending data to inform realistic assumptions about retiree consumption patterns. 

\subsection{Asset Return Data}

The simulation framework requires historical return data for major asset classes, which will be used to drive the Monte Carlo simulations. Consilium maintains access to market return databases from Morningstar, Dimensional, and MSCI, covering New Zealand and international equities, bonds, and cash equivalents. These databases provide total return indices incorporating dividends and interest. We will supplement with publicly available academic datasets to extend the historical coverage where necessary.

The ideal dataset will span at least 50 years to capture multiple complete market cycles, including periods of strong and weak performance, varying inflation, and different interest rate environments. This data will be formatted as a time series of returns for each major asset class, serving as the foundation for generating longer synthetic return paths through resampling techniques. The choice of resampling method is critical, as simple random sampling assumes returns are independent and identically distributed, which may not capture important features and time-dependencies of real market dynamics. In particular, we will explore the block-bootstrapping method used by \textcite{anarkulova2025beyond}, which preserves time-series properties of returns while generating sufficient variation for Monte Carlo analysis. However, without access to their exact dataset, precise replication may not be possible.

\subsection{Retirement Spending Data}

Understanding realistic consumption patterns for New Zealand retirees is essential for developing appropriate spending assumptions in our simulation model. Academic literature on retirement spending has documented important patterns that differ from simple constant-spending assumptions used by many previous studies. In particular, spending tends to decline with age, though not uniformly across all categories, while healthcare costs often increase in later years even as discretionary spending falls.

For the New Zealand context, we will synthesise data from multiple sources from academic research, government statistics, and industry reports. Academic literature, including international studies, will inform our understanding of typical spending trajectories and their variation with wealth levels. However, the most relevant data will come from research conducted specifically on New Zealand retirees. The author has established connections with Ian Parera, who has conducted comprehensive research on retirement consumption in New Zealand through the Retirement Income Interest Group of the New Zealand Society of Actuaries. This research provides detailed insights into how New Zealand retirees actually spend throughout retirement, including the impact of New Zealand Superannuation on spending patterns and the extent of wealth drawdown. By incorporating spending patterns reflecting actual New Zealand retiree behaviour, we can assess whether optimal asset allocation strategies differ meaningfully from those derived under stylised assumptions. This is particularly important given that international literature may not fully capture unique features of New Zealand's retirement income system.

\subsection{Data Challenges}

Several important challenges arise when using historical data to inform lifetime asset allocation strategies.

The first concerns sample size limitations. While long-term return data is available for major asset classes, the effective sample size for meaningful optimisation may be constrained by the enormous number of return paths required for robust Monte Carlo simulation. We need sufficiently long return histories to capture the full range of market behaviour, including rare but important events like sustained bear markets or high inflation periods. \textcite{anarkulova2025beyond} highlight the importance of using bootstrapped returns rather than parametric models to better capture real-world return dynamics, and develop a block-bootstrapping method generating approximately one million synthetic return paths from historical data while preserving important statistical properties. We will study this method in detail and adapt it to our available data if possible.

Another challenge is that policy settings undergo structural changes over time, complicating the use of historical data for future projections. New Zealand's retirement system has undergone significant reforms in just the last two decades. KiwiSaver's introduction in 2007 represents the most substantial change, creating a subsidised retirement savings vehicle unavailable to older cohorts. This may affect both typical wealth levels of future retirees and their spending behaviour, as KiwiSaver balances are held separately from other savings and may be treated differently behaviourally. Changes to New Zealand Superannuation, including adjustments to eligibility age or means testing (should they occur), would fundamentally alter the optimisation problem by changing the guaranteed income floor. Additionally, political support for increasing KiwiSaver is increasing \parencite[e.g.][]{national2025kiwisaver}. We will not attempt to model all possible future reforms. We will assume that current policy settings remain in place as it relates to Superannuation and healthcare, while we may look at testing different KiwiSaver contribution scenarios.

A further issue is the generalisability of spending patterns derived from aggregate data. Retirement spending is inherently heterogeneous, with individual consumption needs varying widely based on health status, housing situation, family circumstances, and lifestyle preferences. Some retirees live primarily on New Zealand Superannuation with minimal discretionary spending, while others maintain much higher consumption levels supported by substantial wealth. Healthcare costs exhibit enormous individual variation, from minimal expenses to substantial ongoing costs. To limit project scope, we will focus on modelling spending patterns reflecting "typical" retirees—those around or below median wealth levels who rely on New Zealand Superannuation and personal savings to fund modest but comfortable retirement lifestyles. We will not model high-net-worth retirees or complex financial situations. However, we will explore scenarios for individuals and couples, as well as low, medium, and high spenders, to assess the robustness of findings across plausible consumption needs.

Any conclusions drawn from historical data are inherently uncertain and based on assumptions that may not hold in the future. This is an inherent limitation that cannot be fully resolved. We must therefore be transparent about our assumptions regarding market returns, spending patterns, and policy settings, clearly communicating both the uncertainty and limitations of our findings.


\section{Planning}

The project breaks into several overlapping workstreams/phases spanning from February through June 2026. Figure \ref{fig:gantt} presents the timeline as a Gantt chart, showing the approximate duration and dependencies between phases. The following sections describe each phase in detail, including key activities and deliverables.


\begin{figure}[ht]
    \centering
    \begin{ganttchart}[
            hgrid,
            vgrid,
            x unit=2cm,
            y unit title=0.6cm,
            y unit chart=1cm, 
            title height=1,
            bar/.style={fill=gray!50},
            bar height=0.6
        ]{1}{5}

        % Title
        \gantttitle{2026}{5} \\
        % Months
        \gantttitle{Feb}{1}
        \gantttitle{Mar}{1}
        \gantttitle{Apr}{1}
        \gantttitle{May}{1}
        \gantttitle{Jun}{1} \\
        % Tasks
        \ganttbar{Phase 1: Literature Review}{1}{1} \\
        \ganttbar{Phase 2: Framework Development}{1}{2} \\
        \ganttbar{Phase 3: Age-Based Glide-Path Optimisation}{2}{3} \\
        \ganttbar{Phase 4: Wealth-Responsive Optimisation}{3}{4} \\
        \ganttbar{Phase 5: Report Writing}{4}{5}
    \end{ganttchart}
    \caption{Project timeline showing the five main phases of work}
    \label{fig:gantt}
\end{figure}


\subsection{Phase 1: Literature Review}

The first phase, to be completed by early February, focuses on establishing the theoretical foundations for the modelling framework and identifying realistic assumptions for New Zealand retirees. There are two main areas of interest that require in-depth review. First, we need to understand the existing literature on lifetime asset allocation, particularly studies that have examined age-based glide-paths and wealth-responsive strategies. This includes foundational works such as \textcite{merton1969lifetime} and more recent contributions like \textcite{anarkulova2025beyond} and \textcite{Makinen2024}. We will critically assess the methodologies used in these studies, their assumptions about spending patterns and risk measures, and the extent to which their findings can be generalised to the New Zealand context. Then we will review the literature on retiree spending behaviour and the goals and concerns of real-world retirees. Because the majority of this corpus is international in scope, we will pay special attention to studies conducted in countries with retirement systems similar to New Zealand's, such as Australia, with a focus on how and why spending patterns differ between countries as a result of policy settings, cultural factors, and economic conditions. We will also closely examine research specifically focused on New Zealand retirees, to extract insights and data that can inform our assumptions about spending patterns and risk concerns. The primary deliverable is a comprehensive literature review section that will form part of the final report, providing both the theoretical foundation for the research and justification for the specific modelling choices made in subsequent phases.


\subsection{Phase 2: Framework Development}

The second phase, to be completed by early March, involves building the core simulation and optimisation infrastructure that will support all subsequent analysis. We will implement a Monte Carlo simulation framework capable of generating wealth trajectories under various lifetime allocation strategies, incorporating the realistic spending rules and withdrawal patterns identified in Phase 1. A critical design decision is ensuring the framework remains agnostic to the source of asset returns, allowing us to test strategies using bootstrapped historical data, parametrically generated returns, or other models. This flexibility is important because different return models may lead to different conclusions about optimal strategies, and we need to assess the robustness of our findings across these alternatives. We will also study and develop risk measures appropriate for lifetime allocation problems, moving beyond traditional risk metrics like variance to include measures that capture the practical concerns of retirees, such as the probability of portfolio depletion at different ages or total consumption. The code will be thoroughly documented to ensure reproducibility and to facilitate potential future extensions by other researchers. The primary deliverable is a working simulation framework with complete documentation, which will serve as the engine for all optimisation work in subsequent phases.


\subsection{Phase 3: Age-Based Glide-Path Optimisation}

The third phase, to be completed by late March, applies the simulation framework to evaluate and optimise age-based glide-path strategies. We will begin by attempting to replicate the findings of \textcite{anarkulova2025beyond} regarding age-based glide-paths, using their methodology as closely as possible. This will also involve implementing their block-bootstrapping method for generating return paths, which we will also seek to adapt for further use in our framework. This replication serves multiple purposes: it validates our framework against published results, helps us understand the sensitivity of their conclusions to methodological choices, and provides a baseline for comparison when we introduce New Zealand-specific assumptions. We will then optimise age-based glide-paths using consumption models based on the New Zealand context and policy settings. This phase will also include a sensitivity analysis to assess how optimal strategies change with different risk measures, return models, and assumptions about spending behaviour. We are particularly interested in whether the conclusions drawn from bootstrapped returns differ meaningfully from those based on parametric return models, as highlighted by \textcite{anarkulova2025beyond}. This is particularly relevant given that much of Consilium's existing advice is based on traditional assumptions modelled with IID asset returns. The primary deliverable is a set of optimised age-based glide-paths along with performance comparisons that can be translated into actionable recommendations for New Zealand retirees and their advisers.


\subsection{Phase 4: Wealth-Responsive Optimisation}

The fourth phase, to be completed by late April if time permits, represents a more ambitious extension of the optimisation framework. Building on the control-matrix approach of \textcite{Makinen2024}, we will develop methods for optimising strategies that respond to both age and wealth levels rather than age alone. This extension is theoretically important because it allows strategies to adapt to an investor's actual financial position rather than assuming all investors of the same age should follow identical allocations. However, implementing this approach requires addressing several limitations in the original Mäkinen method. Their risk measure focuses exclusively on terminal wealth, which we will extend to capture time-dependent concerns about portfolio depletion. Their use of fixed wealth bounds in the control matrix also creates inefficiencies that we will address through adaptive bounds or alternative parameterisation approaches. The optimisation problem in this phase is substantially more complex than in Phase 3, as the parameter space expands significantly when allocation depends on both dimensions.

Given the technical challenges involved, this phase is designated as optional in our planning. If earlier phases take longer than anticipated or if the optimisation proves intractable within the available timeframe, we will focus instead on deepening the analysis from Phase 3. The deliverable, if completed, is a set of wealth-responsive allocation strategies along with performance evaluations that can be compared to the age-based glide-paths developed earlier. These will then be translated into more basic rules-of-thumb that can be communicated effectively to financial advisers and their clients.

\subsection{Phase 5: Analysis and Writing}

The final phase, spanning May and June, focuses on synthesising findings and producing the final research outputs. This involves drawing findings from all previous phases, focusing on identifying robust conclusions that hold across different modelling assumptions, and developing specific recommendations appropriate for the New Zealand context. The recommendations must acknowledge both the strengths and limitations of the analysis, including explicit discussion of the assumptions on which they rest and the circumstances under which they might not apply. We will complete the final report writing during this period, ensuring that all technical work is properly documented and that findings are presented in a manner accessible to both academic and practitioner audiences. Additionally, all code developed during the project will be cleaned up, documented, and published in a public GitHub repository to ensure transparency and reproducibility.

The draft report will be completed by early May to allow time for feedback and revision, with the final version due in June. The primary deliverables are the complete final report, GitHub repo, and presentation materials suitable for sharing with Consilium's advisers and clients.


\subsection{Contingency Planning}

Some contingency measures have been identified to address potential challenges that may arise during the project. As we've mentioned, if Phase 4 proves too complex to complete within the available timeframe, we will redirect effort towards deepening the analysis in Phase 3 rather than rushing to produce incomplete or poorly validated results for the wealth-responsive optimisation. This might include more extensive sensitivity testing, examination of additional risk measures, or analysis of how optimal strategies vary across different investor profiles. The core insights from age-based optimisation are sufficiently valuable that a thorough treatment of Phase 3 alone would constitute a meaningful contribution.

Additionally, if optimisation convergence proves problematic, particularly in Phase 4 where the parameter space is large, we will simplify the strategy space to facilitate convergence. This might involve limiting allocations to 10 percent increments rather than allowing continuous variation, reducing the time granularity of the glide-path, or constraining the wealth dimension of the control matrix to fewer grid points. While this reduces the theoretical optimality of the solution, it ensures we obtain robust and interpretable results rather than solutions that are sensitive to numerical instabilities in the optimisation procedure.


\section{Resources}


\subsection{Data Access}

Market return data will be accessed through Consilium's existing subscriptions to professional financial databases, including Morningstar Direct, Dimensional Fund Advisors' return series, and MSCI indices. These databases provide comprehensive coverage of New Zealand and international equity markets, fixed income securities, and cash equivalents, with return histories extending back several decades in most cases. Consilium has confirmed access to these resources for the duration of the project, with data to be provided in a format suitable for analysis as structured output available via flat files, APIs, and databases. The partnership agreement specifies that the author retains rights to use this data for academic research purposes, including publication, though the raw data itself remains proprietary to the database providers. In addition to Consilium's commercial databases, we will supplement with publicly available academic datasets where appropriate to facilitate replication and comparison of results.

\subsection{Expert Support}

Two external experts have agreed to provide consultation on specific technical aspects of the research, with clearly defined scopes of engagement to ensure expectations are manageable for all parties.

Ian Parera has conducted extensive research on retirement consumption patterns in New Zealand through his work with the Retirement Income Interest Group of the New Zealand Society of Actuaries. He has agreed to provide guidance on realistic spending assumptions for New Zealand retirees, including advice on appropriate consumption trajectories across different age groups and wealth levels, the role of New Zealand Superannuation in meeting spending needs, and how healthcare and other costs evolve through retirement. His support will primarily take the form of two to three structured consultations during Phase 1 (literature review) and early Phase 2 (framework development), where we will discuss his research findings and seek feedback on the spending assumptions we plan to incorporate into the simulation model. This arrangement is informal and advisory in nature, with no expectation of co-authorship or ongoing time commitment beyond these initial consultations, although he has expressed interest in reviewing the final report once completed.

We will also seek academic consultation from experts, particularly as it relates to the optimisation challenges in Phases 3 and 4. The author has a connection with Philipp Wacker at the University of Canterbury, who has offered to provide technical advice on numerical optimisation methods. His support will be particularly valuable for Phase 4, where we attempt to extend the control-matrix approach to incorporate wealth-responsive strategies. He has indicated willingness to review relevant code sections if needed to diagnose optimisation problems, though the primary responsibility for implementation remains with the author. This is an informal advisory arrangement with no expectation of ongoing involvement beyond specific consultations as needed.

\section{Constraints}

\subsection{Data Privacy}

No personal or confidential data will be used in this research. All market return data will be drawn from aggregated indices that are publicly available or accessible through standard financial data providers. Spending data will be synthesised exclusively from published academic research, government statistics, and industry reports. No individual-level retirement account data or personal financial information will be accessed or analysed. This ensures that the research raises no privacy or confidentiality concerns.

It should be noted that Consilium does maintain a large database of client financial information, including detailed records of KiwiSaver balances, contribution histories, and demographic data. However, this data will not be used in the research. The focus is on developing generalisable asset allocation strategies based on aggregate market data and synthesised spending patterns rather than analysing individual client outcomes. This approach avoids any potential privacy issues and ensures that findings are broadly applicable rather than tied to specific client profiles. Consilium could adapt the framework developed here to analyse their own client data in the future, but that would be a separate project requiring appropriate ethical approvals and data governance measures.

Given that proprietary market return data from commercial databases will be used, all results will be presented in aggregate form without disclosing any raw data. The research findings will focus on the performance of asset allocation strategies rather than specific return series, ensuring compliance with data usage agreements and protecting the confidentiality of proprietary information.


\subsection{Scope}

Several important boundaries have been established to ensure the project remains focused and achievable within the available timeframe while still addressing the core research questions.

The project focuses exclusively on strategic asset allocation, the primary driver of investment outcomes \parencite{ibbotsonAssetAllocation}. Analysis will be limited to equities, bonds, and cash, without examining sub-categories within these classes (e.g., large-cap versus small-cap equities, or government versus corporate bonds). This simplification is designed to keep the analysis manageable and focused on the primary drivers of portfolio performance. Our aim is not to recommend a specific portfolio design, but rather to identify general principles for how retirees should allocate to risk across their lifetime.

The research focuses exclusively on the New Zealand context. While international comparisons may provide context or validation against published studies, primary analysis will use New Zealand-specific data and assumptions. This ensures practical value for Consilium and their adviser clients, who work exclusively with New Zealand investors, while providing a natural boundary limiting scope creep. Although the framework could theoretically be adapted to other countries, doing so would require significant additional work beyond this project's scope.

As mentioned in the Data section, the research will focus on modelling "typical" New Zealand retirees rather than attempting to capture the full heterogeneity of retiree circumstances. This means we will primarily consider middle-of-the-road retirees with moderate wealth levels who rely on a combination of New Zealand Superannuation and personal savings to fund a comfortable retirement lifestyle. We will also assume that policy settings with regard to New Zealand Superannuation, KiwiSaver, and healthcare remain largely unchanged over the analysis period. While we may explore different KiwiSaver contribution scenarios, we will not attempt to model the full range of possible future policy reforms. 

While we will explore asset allocation during the accumulation and decumulation phases of retirement, our main focus will be on the decumulation phase, where retirees draw down their savings to fund consumption. This is done to avoid also having to accurately model contribution and labour income patterns during the accumulation phase, which would add significant complexity. We will try to model a range of labour incomes based on typical New Zealand earnings, but the primary focus will be on how retirees should allocate their portfolios once they begin withdrawing funds.


\printbibliography


\end{document}